\chapter{The Sherrington--Kirkpatrick Model at High Temperature: The Talagrand Theory}\label{chap:sg-talagrand}

This chapter formalizes Sections~1.1 through 1.5 of Michel Talagrand's \emph{Spin Glasses: A Challenge for Mathematicians} (Volume~I, Springer, 2010).
The mathematical narrative develops the high-temperature theory of the Sherrington--Kirkpatrick (SK) model:
\begin{enumerate}
\item The decoupled reference system, hypercube factorizations, exact first-moment calculations, and the first-moment upper bound on the free energy (Section~1.2).
\item Guerra's smart path Gaussian interpolation, the derivative formula with a negative square remainder, and Guerra's replica-symmetric (RS) upper bound minimized at the fixed-point equation $q = \mathbb{E}\tanh^2(\beta z \sqrt{q} + h)$ (Section~1.3).
\item R.~Lata\l{}a's high-temperature overlap concentration: the general second-order replica derivative formula along Gaussian interpolations, the monotonicity of the compensated exponential functional, exponential overlap concentration $\nu(\exp s N (R_{1,2} - q)^2) \le (1 - 2s - 4\beta^2)^{-1/2}$ for $\beta < 1/2$, and the Aizenman--Lebowitz--Ruelle (ALR) high-temperature free energy formula $|p_N(\beta, h) - \operatorname{SK}(\beta, h)| \le K/N$ (Section~1.4).
\item Spin decorrelation, asymptotic product measures on finite spin marginals, and high-dimensional random projections (Sections~1.4--1.5).
\end{enumerate}

\medskip\noindent\textbf{Conventions and Notation.}
The spin configuration space is $\Sigma_N = \{-1, +1\}^N$, represented in Lean as functions $\sigma : \mathrm{Fin}\,N \to \mathbb{Z}^\times$ where $\mathbb{Z}^\times = \{-1, +1\}$.
The Hamiltonian with inverse temperature $\beta > 0$ and external field $h \in \mathbb{R}$ is
\[
-H_N(\sigma) = \frac{\beta}{\sqrt{N}} \sum_{1 \le i < j \le N} g_{ij} \sigma_i \sigma_j + h \sum_{i \le N} \sigma_i,
\]
where the disorder $(g_{ij})_{1 \le i < j \le N}$ consists of $\binom{N}{2} = \frac{N(N-1)}{2}$ independent standard Gaussian variables $\mathcal{N}(0, 1)$ indexed by $\mathrm{FinPairs}\,N = \{p \in \mathrm{Fin}\,N \times \mathrm{Fin}\,N \mid p.1 < p.2\}$.
The partition function is $Z_N(\beta) = \sum_\sigma \exp(-H_N(\sigma))$, the expected free energy is $p_N(\beta) = N^{-1}\mathbb{E}\log Z_N(\beta)$, and the overlap of two configurations is $R_{1,2} = R(\sigma^1, \sigma^2) = N^{-1}\sum_{i \le N} \sigma_i^1 \sigma_i^2$.

\section{The SK Hamiltonian and Foundations (Talagrand Section~1.1)}\label{sec:sg-tal:foundations}

In Section~1.1, the SK Hamiltonian is defined on the discrete hypercube $\Sigma_N$.
Its covariance structure, Gaussian comparison properties, concentration of measure, and the thermodynamic limit are foundational.

\begin{definition}[Disorder pairs and Hamiltonian]\label{def:sg-tal:hamiltonian}
  \lean{SpinGlasses.FinPairs, SpinGlasses.card_finPairs, SpinGlasses.cast_card_finPairs}
  \leanok
For $N \in \mathbb{N}$, the type of disordered pairs is $\mathrm{FinPairs}\,N = \{p : \mathrm{Fin}\,N \times \mathrm{Fin}\,N \mid p.1 < p.2\}$.
Its cardinality is $|\mathrm{FinPairs}\,N| = \binom{N}{2} = \frac{N(N-1)}{2}$.
For disorder $g : \mathrm{FinPairs}\,N \to \mathbb{R}$ and configuration $\sigma : \mathrm{Fin}\,N \to \mathbb{Z}^\times$, the Hamiltonian is
\[
-H_N(g, \sigma) = \frac{\beta}{\sqrt{N}} \sum_{p \in \mathrm{FinPairs}\,N} g_p \sigma_{p.1} \sigma_{p.2} + h \sum_{i \le N} \sigma_i.
\]
\end{definition}

\begin{theorem}[Covariance of the SK Hamiltonian]\label{thm:sg-tal:covariance}
For any two configurations $\sigma^1, \sigma^2 \in \Sigma_N$, the centered Hamiltonian has covariance
\[
\mathbb{E} \left[\left(\frac{\beta}{\sqrt{N}} \sum_{i < j} g_{ij} \sigma_i^1 \sigma_j^1\right)\left(\frac{\beta}{\sqrt{N}} \sum_{i < j} g_{ij} \sigma_i^2 \sigma_j^2\right)\right] = \frac{N \beta^2}{4}\left(R(\sigma^1, \sigma^2)^2 - \frac{1}{N}\right).
\]
\end{theorem}
\begin{proof}
By independence of $g_{ij}$, the expectation is $\frac{\beta^2}{N} \sum_{i < j} \sigma_i^1 \sigma_j^1 \sigma_i^2 \sigma_j^2$.
Using $\sum_{i < j} x_i x_j = \frac{1}{2}((\sum_i x_i)^2 - \sum_i x_i^2)$ with $x_i = \sigma_i^1 \sigma_i^2$ (so $x_i^2 = 1$ and $\sum x_i = N R_{1,2}$),
the sum equals $\frac{1}{2}(N^2 R_{1,2}^2 - N) = \frac{N^2}{2}(R_{1,2}^2 - 1/N)$.
Multiplying by $\beta^2 / N$ gives $\frac{N\beta^2}{4}(R_{1,2}^2 - 1/N)$.
\end{proof}

\begin{theorem}[Gaussian comparison and Slepian's lemma]\label{thm:sg-tal:slepian}
  \lean{SpinGlasses.integral_log_sum_exp_le_of_covariance_le}
  \leanok
Let $(X_\sigma), (Y_\sigma)$ be centered Gaussian fields on $\Sigma_N$ such that $\mathbb{E} X_\sigma^2 = \mathbb{E} Y_\sigma^2$ and $\mathbb{E} X_{\sigma^1} X_{\sigma^2} \le \mathbb{E} Y_{\sigma^1} Y_{\sigma^2}$ for all $\sigma^1, \sigma^2$.
Then for any bias $h$,
\[
\mathbb{E} \log \sum_\sigma \exp(Y_\sigma + h_\sigma) \le \mathbb{E} \log \sum_\sigma \exp(X_\sigma + h_\sigma).
\]
\end{theorem}
\begin{proof}
  \leanok
Differentiating along the smooth Slepian path $s \mapsto \cos s X + \sin s Y$ and applying Gaussian integration by parts.
Formalized in \texttt{SpinGlasses/GaussianComparison.lean}.
\end{proof}

\begin{theorem}[Concentration of the free energy]\label{thm:sg-tal:concentration}
  \lean{SpinGlasses.mixedPSpin_free_energy_concentration}
  \leanok
For all $t > 0$, the free energy $F_N = N^{-1}\log Z_N$ satisfies the sub-Gaussian concentration inequality
\[
\mathbb{P}\left(\left|F_N - \mathbb{E} F_N\right| \ge t\right) \le 2 \exp\left(-\frac{N t^2}{2 \beta^2}\right).
\]
\end{theorem}
\begin{proof}
  \leanok
By the Tsirelson--Ibragimov--Sudakov Gaussian concentration inequality, as the gradient of the log-partition function with respect to $g$ has Euclidean norm bounded by $\beta \sqrt{N}$.
Formalized in \texttt{SpinGlasses/Concentration.lean}.
\end{proof}

\begin{theorem}[The thermodynamic limit: Guerra--Toninelli theorem]\label{thm:sg-tal:gt}
  \lean{SpinGlasses.integral_log_sum_exp_add_le, SpinGlasses.sk_tendsto_free_energy}
  \leanok
The sequence of expected free energies $p_N(\beta, h) = N^{-1}\mathbb{E}\log Z_N(\beta, h)$ is superadditive, and the thermodynamic limit exists:
\[
p(\beta, h) = \lim_{N \to \infty} p_N(\beta, h) = \sup_{N \ge 1} p_N(\beta, h).
\]
\end{theorem}
\begin{proof}
  \leanok
By Guerra--Toninelli cavity interpolation splitting a system of size $N_1 + N_2$ into two independent subsystems of sizes $N_1$ and $N_2$.
Formalized in \texttt{SpinGlasses/GuerraToninelli.lean}.
\end{proof}

\section{Decoupled Systems, Hypercube Factorization and First Moments (Talagrand Section~1.2)}\label{sec:sg-tal:decoupled}

Section~1.2 of Talagrand establishes the non-interacting baseline: product factorization over the hypercube, local magnetizations under decoupled fields, exact calculation of $\mathbb{E} Z_N$, the Jensen upper bound, and the temperature derivative of the free energy.

\begin{lemma}[Hypercube product factorization]\label{lem:sg-tal:hypercube-prod}
  \lean{sum_units_int,
        hypercube_sum_prod_eq_prod_add,
        hypercube_sum_prod_erase_eq_prod_add}
  \leanok
For any finite index set $\iota$ and family of functions $a : \iota \to \mathbb{Z}^\times \to \mathbb{R}$,
\[
\sum_{\sigma : \iota \to \mathbb{Z}^\times} \prod_{i \in \iota} a(i, \sigma_i) = \prod_{i \in \iota} (a(i, 1) + a(i, -1)).
\]
For any fixed coordinate $j \in \iota$, the restricted sum satisfies
\[
\sum_{\sigma : \iota \to \mathbb{Z}^\times} \prod_{i \neq j} a(i, \sigma_i) = 2 \prod_{i \neq j} (a(i, 1) + a(i, -1)).
\]
\end{lemma}
\begin{proof}
  \leanok
By induction on the finite set $\iota$ using Fintype embeddings and splitting the sum at $j$.
Formalized in \texttt{ForMathlib/Algebra/BigOperators/HypercubeProduct.lean}.
\end{proof}

\begin{theorem}[Decoupled partition function and local observables]\label{thm:sg-tal:decoupled-partition}
  \lean{Real.continuous_tanh,
        SpinGlasses.decoupled_partition_function_eq,
        SpinGlasses.decoupled_gibbs_spin_eq,
        SpinGlasses.decoupled_gibbs_replica_two_eq,
        SpinGlasses.decoupled_free_energy_eq}
  \leanok
For the decoupled cavity-free Hamiltonian $-H_0(\sigma) = \sum_{i \le N} \sigma_i Y_i$, the partition function factorizes as
\[
Z_0 = 2^N \prod_{i=1}^N \cosh Y_i.
\]
The decoupled Gibbs expectation of a single spin satisfies $\langle \sigma_j \rangle_0 = \tanh Y_j$,
the two-replica overlap satisfies $\langle \sigma_j^1 \sigma_j^2 \rangle_0 = \tanh^2 Y_j$,
and the decoupled free energy is
\[
p_N(0) = \log 2 + \frac{1}{N} \sum_{i=1}^N \mathbb{E} \log\cosh Y_i.
\]
\end{theorem}
\begin{proof}
  \leanok
Apply Lemma~\ref{lem:sg-tal:hypercube-prod} with $a(i, \pm 1) = \exp(\pm Y_i)$ and use $\exp(Y) + \exp(-Y) = 2\cosh Y$ and $\exp(Y) - \exp(-Y) = 2\sinh Y$.
Formalized in \texttt{SpinGlasses/SKDecoupled.lean}.
\end{proof}

\begin{theorem}[Exact first moment of the SK partition function]\label{thm:sg-tal:first-moment}
  \lean{SpinGlasses.integral_exp_mul_gaussianReal_zero_one,
        SpinGlasses.integrable_exp_sum_mul_gaussianReal,
        SpinGlasses.integral_exp_sum_mul_gaussianReal,
        SpinGlasses.sk_first_moment_term_sq,
        SpinGlasses.sum_sk_first_moment_term_sq,
        SpinGlasses.sum_exp_h_mul_spin,
        SpinGlasses.integral_sk_partition_function,
        SpinGlasses.integral_sk_partition_function_eq}
  \leanok
For the SK model with Gaussian disorder $g_p \sim \mathcal{N}(0, 1)$ on pairs $1 \le i < j \le N$ and constant external field $h \in \mathbb{R}$,
the exact first moment of the partition function is
\[
\mathbb{E} Z_N = (2\cosh h)^N \exp\left(N \frac{\beta^2}{4}\left(1 - \frac{1}{N}\right)\right),
\]
which yields the exact normalized log-moment (Talagrand Exercise~1.2.3, Eq.~(1.25)):
\[
\frac{1}{N} \log \mathbb{E} Z_N = \frac{\beta^2}{4}\left(1 - \frac{1}{N}\right) + \log 2 + \log\cosh h.
\]
\end{theorem}
\begin{proof}
  \leanok
For each configuration $\sigma$, the Gaussian linear combination $\sum_{i<j} \frac{\beta}{\sqrt{N}} g_{ij} \sigma_i \sigma_j$ has variance $\sum_{i<j} \frac{\beta^2}{N} = \binom{N}{2}\frac{\beta^2}{N} = N \frac{\beta^2}{4}(1 - \frac{1}{N})$, independent of $\sigma$.
Integrating out the Gaussian disorder factors the sum into $(2\cosh h)^N \exp(N \frac{\beta^2}{4}(1 - 1/N))$.
Formalized in \texttt{SpinGlasses/SKFirstMoment.lean}.
\end{proof}

\begin{theorem}[First-moment Jensen upper bound on the free energy]\label{thm:sg-tal:first-moment-bound}
  \lean{MeasureTheory.integral_log_le_log_integral,
        SpinGlasses.sk_free_energy_le_first_moment}
  \leanok
By Jensen's inequality for the concave logarithm under the disorder probability measure, the expected free energy satisfies (Eqs.~(1.24) and (1.26)):
\[
p_N(\beta, h) \le \frac{1}{N} \log \mathbb{E} Z_N = \frac{\beta^2}{4}\left(1 - \frac{1}{N}\right) + \log 2 + \log\cosh h.
\]
\end{theorem}
\begin{proof}
  \leanok
$p_N(\beta, h) = N^{-1}\mathbb{E}\log Z_N \le N^{-1}\log \mathbb{E} Z_N$, and apply Theorem~\ref{thm:sg-tal:first-moment}.
Formalized in \texttt{SpinGlasses/SKFirstMoment.lean}.
\end{proof}

\begin{theorem}[Temperature derivative of the expected free energy]\label{thm:sg-tal:beta-derivative}
  \lean{SpinGlasses.hasDerivAt_log_sum_exp_beta,
        SpinGlasses.measurable_disorder_u,
        SpinGlasses.memLp_finPairs_entry,
        SpinGlasses.integrable_finPairs_entry_abs,
        SpinGlasses.hasDerivAt_sk_free_energy_beta}
  \leanok
For fixed disorder $g$, the map $\beta \mapsto \log Z_N(\beta)$ has derivative equal to the Gibbs average $\langle u_\sigma \rangle$.
Under the product Gaussian disorder on pairs $i < j$, dominated differentiation under the integral yields Lemma~1.3.11 (Eq.~(1.84)):
\[
\frac{\partial p_N}{\partial \beta} = \frac{1}{N} \mathbb{E} \langle u_\sigma \rangle = \frac{\beta}{2}\left(1 - \mathbb{E}\langle R_{1,2}^2 \rangle\right).
\]
\end{theorem}
\begin{proof}
  \leanok
Differentiating $s \mapsto \log \sum_\sigma \exp(s u_\sigma(g) + C_\sigma)$ yields the Gibbs average.
Dominated differentiation under the integral applies via the integrable bound $B(g) = (\sqrt{N})^{-1}\sum_{i<j} |g_{ij}|$.
Applying Gaussian integration by parts to $\mathbb{E}\langle u_\sigma \rangle$ evaluates the expectation to $\frac{N\beta}{2}(1 - \mathbb{E}\langle R_{1,2}^2 \rangle)$.
Formalized in \texttt{SpinGlasses/SKFirstMoment.lean}.
\end{proof}

\section{Guerra's Smart Path Interpolation and the RS Upper Bound (Talagrand Section~1.3)}\label{sec:sg-tal:guerra-rs}

Section~1.3 introduces Guerra's smart path interpolation between the fully interacting SK model and decoupled Gaussian spins, leading to the replica-symmetric upper bound.

\begin{remark}[Deviations from the book in Section~1.3]\label{rem:sg-tal:dev-13}
\begin{itemize}
\item \emph{Parametrization of the path.} The book interpolates with $\sqrt t\,X+\sqrt{1-t}\,Y$, which is not differentiable at the endpoints. The Lean development uses the trigonometric path $s\mapsto\cos s\,X+\sin s\,Y+h$, $s\in[0,\pi/2]$, with $s=0$ the SK model and $s=\pi/2$ the decoupled model (so $t=\cos^2 s$). All derivative formulas below carry the factor $\sin s\cos s$ instead of the book's $dt$.
\item \emph{Joint Gaussian space.} The pair disorder $g$ and the site disorder $z$ are realized on the product space $(\mathrm{FinPairs}\,N\to\mathbb R)\times(\mathrm{Fin}\,N\to\mathbb R)$ with the product measure; joint Gaussianity of $(X_\sigma,Y_\sigma)_\sigma$ is proved from independence (\texttt{HasGaussianLaw.sumElim\_prod\_of\_independent}).
\item \emph{Scope of the fixed-point statement.} Only the existence of a solution $q\in[0,1]$ of $q=\mathbb{E}\tanh^2(\beta z\sqrt q+h)$ is formalized (intermediate value theorem). The minimization property of Guerra's bound at $q$, the bound $\psi'\le\beta^2$ and uniqueness for $\beta<1$ (Proposition~1.3.8) are not. Every later theorem takes the fixed-point equation as a hypothesis. Guerra's bound itself holds for every $q\ge0$.
\end{itemize}
\end{remark}

\begin{definition}[Guerra's smart path Hamiltonian]\label{def:sg-tal:smart-path}
For $t \in [0, 1]$, independent standard Gaussians $(g_{ij})_{i<j}$ and $(z_i)_{i \le N}$, parameter $q \in [0, 1]$, and bias $h$, the smart path Hamiltonian is
\[
-H_t(\sigma) = \sqrt{t} \frac{\beta}{\sqrt{N}} \sum_{1 \le i < j \le N} g_{ij} \sigma_i \sigma_j + \sqrt{1-t} \beta \sqrt{q} \sum_{i \le N} z_i \sigma_i + h \sum_{i \le N} \sigma_i.
\]
The interpolating partition function is $Z_t = \sum_\sigma \exp(-H_t(\sigma))$ and $\varphi(t) = N^{-1}\mathbb{E}\log Z_t$.
\end{definition}

\begin{lemma}[Covariance of the decoupled Gaussian cavity field]\label{lem:sg-tal:decoupled-covariance}
  \lean{SpinGlasses.integral_decoupled_gaussian_field_mul}
  \leanok
For arbitrary finite real vectors $\sigma, \tau$ and independent standard Gaussian variables
$(z_i)$, the cavity-field part of Guerra's path satisfies
\[
\mathbb{E}\left[\left(\beta\sqrt{q}\sum_i z_i\sigma_i\right)
  \left(\beta\sqrt{q}\sum_i z_i\tau_i\right)\right]
= \beta^2 q \sum_i \sigma_i\tau_i,
\qquad q \geq 0.
\]
\end{lemma}
\begin{proof}
  \leanok
This is the covariance of two linear forms in an iid standard Gaussian vector.
The Lean statement is deliberately for arbitrary finite real vectors, not only hypercube spins,
so it is reusable for relaxed and spherical spin models.  Formalized in
\texttt{SpinGlasses/SKReplicaSymmetricBound.lean}.
\end{proof}

\begin{lemma}[Pair-indexed Gaussian covariance]\label{lem:sg-tal:pair-covariance}
  \lean{SpinGlasses.integral_normalized_gaussian_linear_form_mul,
        SpinGlasses.integral_sk_pair_field_mul,
        SpinGlasses.sum_finPairs_mul_eq_sub_sum_sq_div_two,
        SpinGlasses.sum_sk_pair_coefficient_eq_overlap,
        SpinGlasses.integral_sk_pair_field_mul_eq_overlap}
  \leanok
For pair-indexed iid standard Gaussian disorder and two spin configurations, integrating out the
disorder gives
\[
\mathbb{E}\left[\left(\frac{\beta}{\sqrt{N}}\sum_{i<j}
  g_{ij}\sigma_i\sigma_j\right)
  \left(\frac{\beta}{\sqrt{N}}\sum_{i<j}
  g_{ij}\tau_i\tau_j\right)\right]
= \frac{\beta^2}{N}\sum_{i<j}\sigma_i\sigma_j\tau_i\tau_j.
\]
\end{lemma}
\begin{proof}
  \leanok
The general Lean theorem is the covariance of two normalized linear forms in an iid Gaussian
vector; the SK theorem instantiates its coefficients with the pair products of the two spin
configurations.  A separate finite combinatorial lemma rewrites the unordered pair sum as half
the square of the overlap minus its diagonal, yielding the displayed SK covariance exactly.
Formalized in \texttt{SpinGlasses/SKReplicaSymmetricBound.lean}.
\end{proof}

\begin{lemma}[Algebraic assembly of the smart-path covariance difference]
  \label{lem:sg-tal:path-covariance-algebra}
  \lean{SpinGlasses.smart_path_half_covariance_difference_eq,
        SpinGlasses.sk_smart_path_half_covariance_difference_eq}
  \leanok
For Ising configurations, the two concrete Gaussian fields have covariances
\[
C = \frac{N\beta^2}{2}\left(R(\sigma,\tau)^2 - \frac1N\right),
\qquad
D = \beta^2 q\, \sigma\mathbin{\cdot}\tau.
\]
Then their normalized half-difference is exactly
\[
\frac{C-D}{2N} = \frac{\beta^2}{4}
  \left((R(\sigma,\tau)-q)^2-q^2-\frac1N\right).
\]
\end{lemma}
\begin{proof}
  \leanok
The concrete theorem first invokes the already-formalized pair and cavity covariance calculations.
It then reuses the general algebraic bridge, substituting
$\sigma\mathbin{\cdot}\tau=N R(\sigma,\tau)$ and completing the square.
\end{proof}

\begin{lemma}[Normalized trigonometric smart-path derivative]
  \label{lem:sg-tal:trig-smart-path-derivative}
  \lean{SpinGlasses.hasDerivAt_normalized_log_sum_exp_trig_of_covariance_difference}
  \leanok
Let $X_\sigma$ and $Y_\sigma$ be centered, jointly Gaussian, independent fields on a finite
configuration space, and let $r(\sigma,\tau)$ satisfy $r(\sigma,\sigma)=1$.  Suppose
\[
\frac1N\frac{\operatorname{Cov}(X_\sigma,X_\tau)-
  \operatorname{Cov}(Y_\sigma,Y_\tau)}2
= \frac{\beta^2}{4}\left((r(\sigma,\tau)-q)^2-q^2-\frac1N\right).
\]
For the smooth path $\cos s\,X+\sin s\,Y+h$, the normalized log-partition function has
derivative
\[
\frac{\beta^2}{2}\sin s\cos s\,
\mathbb E\left\langle (r(\sigma^1,\sigma^2)-q)^2-(1-q)^2\right\rangle_s.
\]
\end{lemma}
\begin{proof}
  \leanok
The formal theorem invokes the existing Gaussian integration-by-parts derivative theorem for the
trigonometric path and substitutes the covariance hypothesis pointwise.  It deliberately applies
to arbitrary finite Gaussian fields, so the SK covariance calculation is only one specialization.
This is not yet Eq.~(1.65) for the square-root parameterization; constructing the concrete SK
and cavity fields as a jointly Gaussian pair and matching their endpoints remains part of
Milestone T2.
\end{proof}

\begin{lemma}[Trigonometric replica-symmetric endpoint estimate]
  \label{lem:sg-tal:trig-rs-endpoint}
  \lean{SpinGlasses.le_endpoint_of_trig_square_deriv}
  \leanok
If a differentiable function satisfies
\[
\phi'(s)=\frac{\beta^2}{2}\sin s\cos s
  \left(A(s)-(1-q)^2\right), \qquad A(s)\geq 0,
\]
then
\[
\phi(0)\leq \phi(\pi/2)+\frac{\beta^2}{4}(1-q)^2.
\]
\end{lemma}
\begin{proof}
  \leanok
The formal proof applies Mathlib's monotonicity theorem to
$\phi(s)+\frac{\beta^2}{4}(1-q)^2\sin^2 s$ on $[0,\pi/2]$.
The derivative is $\frac{\beta^2}{2}\sin s\cos s\,A(s)$, which is nonnegative.
\end{proof}

\begin{lemma}[Smart path covariance difference]\label{lem:sg-tal:path-covariance}
Let $u_\sigma = \frac{\beta}{\sqrt{N}}\sum_{i<j} g_{ij} \sigma_i \sigma_j$ and $v_\sigma = \beta \sqrt{q}\sum_i z_i \sigma_i$.
The covariance difference matrix $U(\sigma^1, \sigma^2) = \mathbb{E} u_{\sigma^1} u_{\sigma^2} - \mathbb{E} v_{\sigma^1} v_{\sigma^2}$ satisfies (Eq.~(1.64)):
\[
\frac{1}{N} U(\sigma^1, \sigma^2) = \frac{\beta^2}{4}\left(R_{1,2}^2 - \frac{1}{N}\right) - \frac{\beta^2}{2} q R_{1,2} = \frac{\beta^2}{4}\left((R_{1,2} - q)^2 - q^2 - \frac{1}{N}\right).
\]
On the diagonal $\sigma^1 = \sigma^2$ (where $R_{1,1} = 1$),
\[
\frac{1}{N} U(\sigma, \sigma) = \frac{\beta^2}{4}\left((1 - q)^2 - q^2 - \frac{1}{N}\right).
\]
\end{lemma}
\begin{proof}
Compute $\mathbb{E} u_{\sigma^1} u_{\sigma^2} = \frac{N\beta^2}{4}(R_{1,2}^2 - 1/N)$ and $\mathbb{E} v_{\sigma^1} v_{\sigma^2} = N \beta^2 q R_{1,2}$.
Subtracting and completing the square yields the identity.
\end{proof}

\begin{theorem}[Derivative along the smart path]\label{thm:sg-tal:smart-path-deriv}
The derivative of $\varphi(t) = N^{-1}\mathbb{E}\log Z_t$ along the smart path satisfies (Eq.~(1.65)):
\[
\varphi'(t) = \frac{1}{2N} \left(\mathbb{E}\langle U(\sigma, \sigma) \rangle_t - \mathbb{E}\langle U(\sigma^1, \sigma^2) \rangle_t\right) = -\frac{\beta^2}{4} \mathbb{E}\langle (R_{1,2} - q)^2 \rangle_t + \frac{\beta^2}{4}(1 - q)^2.
\]
In particular, $\varphi'(t) \le \frac{\beta^2}{4}(1 - q)^2$ for all $t \in [0, 1]$.
\end{theorem}
\begin{proof}
By Gaussian integration by parts on the interpolating Gaussian field, the derivative is $\frac{1}{2N}(\mathbb{E}\langle U(\sigma, \sigma) \rangle_t - \mathbb{E}\langle U(\sigma^1, \sigma^2) \rangle_t)$.
Substituting the expressions from Lemma~\ref{lem:sg-tal:path-covariance}, the $-q^2 - 1/N$ terms cancel, leaving $-\frac{\beta^2}{4}\mathbb{E}\langle (R_{1,2} - q)^2 \rangle_t + \frac{\beta^2}{4}(1-q)^2$.
\end{proof}

\begin{theorem}[Guerra's replica-symmetric upper bound]\label{thm:sg-tal:guerra-rs-bound}
  \lean{ProbabilityTheory.HasGaussianLaw.sumElim_prod_of_independent,
        ProbabilityTheory.hasGaussianLaw_dotProduct_family,
        SpinGlasses.hasGaussianLaw_const_mul_sum_family,
        SpinGlasses.sk_pair_sum_eq_dotProduct,
        SpinGlasses.exists_sk_smart_path_fields,
        SpinGlasses.integrable_log_cosh_affine_gaussianReal,
        SpinGlasses.sk_replica_symmetric_upper_bound}
  \leanok
For all $N\ge1$, $\beta, h \in \mathbb{R}$, and any parameter $q \ge 0$ (Theorem~1.3.7):
\[
p_N(\beta, h) \le \log 2 + \mathbb{E} \log\cosh(\beta z \sqrt{q} + h) + \frac{\beta^2}{4}(1 - q)^2,
\]
where $z \sim \mathcal{N}(0, 1)$.
\end{theorem}
\begin{proof}
\leanok
In Lean the interpolation runs along the trigonometric path $s\mapsto\cos s\,X+\sin s\,Y+h$, $s\in[0,\pi/2]$, with $X$ the SK field and $Y$ the cavity field on the product space of the two independent disorders (jointly Gaussian by
\texttt{HasGaussianLaw.sumElim\_prod\_of\_independent}, which allows different index types for the two fields). By Lemma~\ref{lem:sg-tal:trig-smart-path-derivative} the derivative of
$\varphi(s)=N^{-1}\mathbb E\log Z_s$ is $\frac{\beta^2}{2}\sin s\cos s\,(A(s)-(1-q)^2)$ with $A(s)=\mathbb E\langle (R_{1,2}-q)^2\rangle_s\ge0$ (the expectation of a square under a product of two probability measures), so Lemma~\ref{lem:sg-tal:trig-rs-endpoint} gives
$\varphi(0)\le\varphi(\pi/2)+\frac{\beta^2}{4}(1-q)^2$. At $s=0$ one has the SK free energy $p_N(\beta,h)$; at $s=\pi/2$ the system is decoupled with field $Y_i=\beta z_i\sqrt q+h$, so $\varphi(\pi/2)=\log 2+\mathbb E\log\cosh(\beta z\sqrt q+h)$ by Theorem~\ref{thm:sg-tal:decoupled-partition}.
Formalized in \texttt{SpinGlasses/SKGuerraBound.lean} (the free energy profile along the path is \texttt{SKSmartPathProfile.lean}).
\end{proof}

\begin{theorem}[The self-consistency equation and contractivity]\label{thm:sg-tal:fixed-point}
  \lean{SpinGlasses.continuous_fixedPointMap,
        SpinGlasses.fixedPointMap_mem_Icc,
        SpinGlasses.exists_fixedPoint}
  \leanok
\emph{Formalized:} for all $\beta,h\in\mathbb R$ the map $\psi(x)=\mathbb E\tanh^2(\beta z\sqrt x+h)$ is continuous with values in $[0,1]$, hence (intermediate value theorem) the self-consistency equation has a solution $q\in[0,1]$.
\emph{Not formalized:} the minimization property, the bound $\psi'\le\beta^2$ and the uniqueness of the fixed point for $\beta<1$ (they need differentiation of $x\mapsto\mathbb E f(\beta z\sqrt x+h)$ at $x=0$); the later theorems take the fixed-point equation as a hypothesis.
The right-hand side of Guerra's bound is minimized at the fixed point of the self-consistency equation (Eq.~(1.74)):
\[
q = \mathbb{E}\tanh^2(\beta z \sqrt{q} + h).
\]
Defining the replica-symmetric free energy $\operatorname{SK}(\beta, h) = \log 2 + \mathbb{E}\log\cosh(\beta z \sqrt{q} + h) + \frac{\beta^2}{4}(1-q)^2$ at this fixed point,
\[
p_N(\beta, h) \le \operatorname{SK}(\beta, h).
\]
Furthermore, for $\beta < 1$, the map $\psi(q) = \mathbb{E}\tanh^2(\beta z \sqrt{q} + h)$ satisfies $\psi'(q) \le \beta^2 < 1$ (Proposition~1.3.8), ensuring a unique fixed point in $[0, 1]$.
\end{theorem}
\begin{proof}
Differentiating $f(q) = \mathbb{E}\log\cosh(\beta z \sqrt{q} + h) + \frac{\beta^2}{4}(1-q)^2$ with respect to $q$:
using Gaussian integration by parts on $z$, $\frac{d}{dq}\mathbb{E}\log\cosh(\beta z \sqrt{q} + h) = \frac{\beta^2}{2}\mathbb{E}\tanh^2(\beta z \sqrt{q} + h)$.
Setting $f'(q) = \frac{\beta^2}{2}(\mathbb{E}\tanh^2(\beta z \sqrt{q} + h) - q) = 0$ yields the fixed-point equation.
For contractivity, Gaussian integration by parts shows $\psi'(q) = \beta^2 \mathbb{E}(1 - \tanh^2(\beta z \sqrt{q} + h))^2 \le \beta^2 < 1$.
\end{proof}

\section{Lata\l{}a's Overlap Concentration and the ALR Formula (Talagrand Section~1.4)}\label{sec:sg-tal:latala-alr}

Section~1.4 proves R.~Lata\l{}a's exponential overlap concentration theorem and establishes the Aizenman--Lebowitz--Ruelle (ALR) high-temperature free energy formula.

\begin{remark}[Deviations from the book in Section~1.4]\label{rem:sg-tal:dev-14}
Every statement below is proved in Lean, but several proofs differ from Talagrand's, and a few statements are weaker or reorganized.
\begin{itemize}
\item \emph{Lemma~1.4.2: trigonometric form, no $\tfrac12$.} The derivative is taken in the trigonometric parameter $s$, with $U(\sigma,\tau)=\mathbb{E}Y_\sigma Y_\tau-\mathbb{E}X_\sigma X_\tau$ (no factor $1/2$) and prefactor $\sin s\cos s$. Via $t=\cos^2 s$ this is the book's Eq.~(1.90), with the opposite sign because $s=0$ is the SK end. Differentiation under $\mathbb{E}$ is justified explicitly by domination. The lemma is also proved for an observable that depends differentiably on $s$, which is what the compensated functional requires.
\item \emph{Eq.~(1.97) and Corollary~1.4.6 are not separate statements.} The case $n=2$ is expanded inside the proof of Lemma~\ref{lem:sg-tal:replica-decoupling}.
\item \emph{Lemma~1.4.5 by Chebyshev's inequality.} The book expands the exponential in a series, applies H\"older's inequality and uses replica symmetry. Lean applies Chebyshev's association inequality $\int a\int b\le\int ab$ for similarly ordered functions to $a=A$, $b=e^{uA}$, separately for each realization of the disorder. Lemma~1.4.5 is not stated on its own as an inequality between $\nu_t$-averages.
\item \emph{Corollary~1.4.7: direction and generality.} The compensated functional is nondecreasing in the trigonometric parameter $s$, equivalently non-increasing in the book's $t=\cos^2 s$. It is proved abstractly for any nonnegative symmetric overlap function $A$ with constant diagonal and $U=-\kappa(A-c)$, and then specialized to $A=(R_{1,2}-q)^2$.
\item \emph{Proposition~1.4.8.} The one-dimensional inequality $(\cosh u+q\sinh u)e^{-qu}\le e^{u^2/2}$ is obtained from Hoeffding's lemma (Mathlib) applied to an explicit two-point law, not from a direct calculus argument. The randomization $u\sim\mathcal N(0,2\lambda/N)$ is carried out as a Fubini exchange with a standard Gaussian, valid because the overlap is bounded.
\item \emph{Theorem~1.4.1 and the moment bounds.} The exponential bound is proved for the SK endpoint. The uniform bound $\nu_t((R_{1,2}-q)^2)\le K/N$ (Eqs.~(1.89), (1.103)) is derived from the compensated inequality and $e^x\ge1+x$ with an explicit constant; Eq.~(1.88) for all even moments is not formalized.
\item \emph{Theorem~1.4.10 by monotonicity.} Instead of integrating $\varphi'$ over $[0,\pi/2]$, Lean shows that $\varphi+c\sin^2$ is monotone for two choices of $c$ (using $0\le A\le K/N$), which gives the same two-sided $O(1/N)$ bound without any integrability check. The convergence $p_N\to\operatorname{SK}(\beta,h)$ is derived from this bound directly for the pair disorder with field, not from the thermodynamic limit of Theorem~1.3.9.
\item \emph{Hypotheses.} $q$ is a hypothesis (fixed-point equation), and the book's $\beta<1/2$ is stated as $\beta^2<1/4$.
\end{itemize}
\end{remark}

\begin{lemma}[General replica derivative under Gaussian interpolation]\label{lem:sg-tal:replica-deriv}
  \lean{SpinGlasses.hasDerivAt_integral_pi_count_tilted_path,
        SpinGlasses.hasDerivAt_integral_pi_count_tilted_path_moving,
        SpinGlasses.aemeasurable_gibbsWeight,
        SpinGlasses.aemeasurable_integral_pi_count_tilted,
        SpinGlasses.hasDerivAt_integral_replica_path_moving,
        SpinGlasses.hasDerivAt_integral_replica_path,
        SpinGlasses.integrable_integral_pi_count_tilted_mul_eval,
        SpinGlasses.integrable_integral_pi_count_tilted_mul_integral,
        SpinGlasses.integrable_integral_pi_count_tilted_det,
        SpinGlasses.integral_pi_option_count_tilted_marginal,
        SpinGlasses.integral_replica_trig_key,
        SpinGlasses.hasDerivAt_integral_replica_trig,
        SpinGlasses.integrable_replica_derivative_term,
        SpinGlasses.hasDerivAt_integral_replica_trig_moving}
  \leanok
Let $\alpha$ be a finite nonempty configuration space, $R$ a finite set of $n$ replicas, $(\Omega,P)$ a probability space, and $X, Y : \alpha \to \Omega \to \mathbb{R}$ measurable with $(X_\sigma, Y_\sigma)_\sigma$ jointly Gaussian, centered, and $\mathbb{E} X_\sigma Y_\tau = 0$ for all $\sigma,\tau$.
Put $U(\sigma,\tau) = \mathbb{E} Y_\sigma Y_\tau - \mathbb{E} X_\sigma X_\tau$, let $h : \alpha \to \mathbb{R}$, and let $\langle\cdot\rangle_{s,\omega}$ be the average over $n$ independent replicas of the Gibbs measure with weights $\exp(\cos s\, X_\sigma(\omega) + \sin s\, Y_\sigma(\omega) + h_\sigma)$.
For every observable $\Phi : \alpha^R \to \mathbb{R}$ (Lemma~1.4.2, trigonometric form),
\[
\frac{d}{ds}\mathbb{E}\langle\Phi\rangle_s
= \sin s\cos s\;\mathbb{E}\Big\langle \Phi\Big(\sum_{\ell,\ell'\le n} U(\sigma^\ell,\sigma^{\ell'}) - 2n\sum_{\ell\le n} U(\sigma^\ell,\sigma^{n+1}) - n\,U(\sigma^{n+1},\sigma^{n+1}) + n(n+1)\,U(\sigma^{n+1},\sigma^{n+2})\Big)\Big\rangle_s,
\]
where the right-hand side is a Gibbs average over $n+2$ replicas indexed by $\mathrm{Option}(\mathrm{Option}\,R)$.
More generally, if $\Phi_s$ depends differentiably on $s$ with $|\Phi_s|, |\partial_s\Phi_s|$ bounded uniformly in $s$, then $\frac{d}{ds}\mathbb{E}\langle\Phi_s\rangle_s$ is the displayed expression for $\Phi = \Phi_s$ plus $\mathbb{E}\langle\partial_s\Phi_s\rangle_s$.
\end{lemma}
\begin{proof}
\leanok
Differentiate the replica average $\prod_\ell w(\sigma^\ell)/Z^n$ along the path (finite-dimensional calculus, \texttt{fderiv\_integral\_pi\_count\_tilted}) and pass the derivative inside $\mathbb{E}$ by dominated differentiation (the path derivative is dominated by $\sum_\sigma |X_\sigma|+|Y_\sigma|$, which is integrable for a Gaussian field).
This gives $\mathbb{E}\langle\Phi(\sum_\ell u'(\sigma^\ell) - n\langle u'\rangle)\rangle$ with $u'=-\sin s\,X+\cos s\,Y$.
Two Gaussian integrations by parts (\texttt{integral\_gibbs\_replica\_average\_gaussian\_option}, applied to $X$ and to $Y$) turn this into the displayed covariance combination, with $\mathbb{E} X Y = 0$ killing the cross terms; the extra replicas $n+1,n+2$ appear exactly as in Panchenko's Eq.~(1.29), and a marginalization lemma removes the replicas on which the observable does not depend.
In the parametrization $t = \cos^2 s$ (with $t=1$ the SK end) this is Talagrand's Eq.~(1.90); Eq.~(1.97) for $n=2$ is the specialization used in the proof of Lemma~\ref{lem:sg-tal:replica-decoupling}.
Formalized in \texttt{SpinGlasses/GaussianInterpolationReplica.lean}.
\end{proof}

\begin{lemma}[Replica decoupling]\label{lem:sg-tal:replica-decoupling}
  \lean{MeasureTheory.integral_mul_integral_le_integral_mul,
        MeasureTheory.integral_pi_fin_two,
        SpinGlasses.integral_comb_four,
        SpinGlasses.replica_decoupling_comb_nonneg}
  \leanok
Let $\mu$ be a Gibbs measure on a finite set $\alpha$, $A : \alpha\times\alpha\to[0,\infty)$ symmetric with constant diagonal, $U(\sigma,\tau) = -\kappa (A(\sigma,\tau) - c)$ with $\kappa\ge0$, and $u\ge 0$.
For $\Phi(\sigma^1,\sigma^2) = \exp(u A(\sigma^1,\sigma^2))$ the combination of Lemma~\ref{lem:sg-tal:replica-deriv} (for $n=2$) and the moving-observable term for $u'=8\kappa\sin s\cos s$ is nonnegative:
\[
\Big\langle \Phi\,\big(\textstyle\sum U \ldots\big)\Big\rangle_{\mu^{\otimes4}} + 8\kappa\,\langle A\,\Phi\rangle_{\mu^{\otimes2}} \ge 0 .
\]
(For $A=(R_{1,2}-q)^2$ this contains Lemma~1.4.5, $\nu_t((R_{3,4}-q)^2 e^{\lambda N (R_{1,2}-q)^2}) \le \nu_t((R_{1,2}-q)^2 e^{\lambda N (R_{1,2}-q)^2})$, and Corollary~1.4.6.)
\end{lemma}
\begin{proof}
\leanok
\emph{Differs from the book:} Talagrand proves Lemma~1.4.5 by expanding the exponential, H\"older's inequality and replica symmetry.
The Lean proof is shorter and works pointwise in the disorder: since $U$ is $-\kappa(A-c)$ with $A$ symmetric and of constant diagonal, the four-replica combination collapses to $-\kappa(2A_{12}-4A_{13}-4A_{23}+6A_{34})$ (the constants $c$ and the diagonal values cancel because the coefficients sum to zero).
The terms with $A_{13}, A_{23}$ are integrals of nonnegative functions; the term with $A_{34}$ factors by independence of the replicas $3,4$ from $1,2$ into $\langle\Phi\rangle\langle A\rangle$, and Chebyshev's association inequality (for two similarly ordered functions, $\int a\int b\le\int ab$; here $a=A$ and $b=e^{uA}$ with $u\ge0$) gives $\langle A\rangle\langle e^{uA}\rangle\le\langle A e^{uA}\rangle$.
Hence the combination is at least $\kappa(6\langle\Phi A\rangle - 6\langle\Phi\rangle\langle A\rangle)\ge0$.
\emph{Formalization note.} Lemma~1.4.5 is not stated on its own as an inequality between $\nu_t$-averages; only the combination above (Lemma~1.4.2 for $n=2$ plus the moving-observable term), which is what Corollary~\ref{cor:sg-tal:compensated-mono} uses, is a Lean statement.
Formalized in \texttt{SpinGlasses/SKLatalaCompensated.lean} (Chebyshev's inequality in \texttt{ForMathlib/MeasureTheory/Integral/ChebyshevAssociation.lean}).
\end{proof}

\begin{corollary}[Compensated functional monotonicity]\label{cor:sg-tal:compensated-mono}
  \lean{SpinGlasses.exists_sk_smart_path_fields,
        SpinGlasses.monotoneOn_compensated_functional,
        SpinGlasses.sk_compensated_path_le_decoupled,
        SpinGlasses.sk_compensated_le_decoupled}
  \leanok
Abstract form: under the hypotheses of Lemma~\ref{lem:sg-tal:replica-deriv} with covariance difference $U=-\kappa(A-c)$ as in Lemma~\ref{lem:sg-tal:replica-decoupling} and $u_0\ge4\kappa$, the function
$s\mapsto\mathbb{E}\langle\exp((u_0-4\kappa\cos^2 s)A(\sigma^1,\sigma^2))\rangle_s$ is nondecreasing on $[0,\pi/2]$.
SK form (Corollary~1.4.7 in the endpoint version used for Theorem~1.4.1): for $N\ge1$, $q\ge0$, $s\ge0$,
\[
\nu_{\mathrm{SK}}\big(\exp sN(R_{1,2}-q)^2\big) \le \nu_{\mathrm{dec}}\big(\exp (s+2\beta^2) N (R_{1,2}-q)^2\big),
\]
where $\nu_{\mathrm{dec}}$ is the average for the decoupled model with fields $\beta\sqrt q\, z_i + h$.
\end{corollary}
\begin{proof}
\leanok
By Lemma~\ref{lem:sg-tal:replica-deriv} (moving-observable form, $\Phi_s=\exp((u_0-4\kappa\cos^2s)A)$, $\partial_s u = 8\kappa\sin s\cos s$) the derivative of the functional is $\sin s\cos s$ times the nonnegative quantity of Lemma~\ref{lem:sg-tal:replica-decoupling} (note $u_0-4\kappa\cos^2 s\ge0$), so it is nonnegative on $[0,\pi/2]$.
For the SK model take $X$ the pair field and $Y$ the cavity field (\texttt{exists\_sk\_smart\_path\_fields} packages their Gaussianity and Eq.~(1.64)), $A=(R-q)^2$, $\kappa=N\beta^2/2$, $c=q^2+1/N$, $u_0=N(s+2\beta^2)$, so $u_0-4\kappa=Ns$; the endpoints $s=0$ and $s=\pi/2$ are the SK and decoupled models.
In Talagrand's parametrization $t=\cos^2 s$ the functional $\nu_t(\exp((\lambda-2t\beta^2)N(R_{1,2}-q)^2))$ is non-increasing in $t$, with $\lambda=s+2\beta^2$.
The path version compares the smart-path energy at any time $t\in[0,\pi/2]$ with the decoupled endpoint, which gives the uniform moment bounds of Theorem~\ref{thm:sg-tal:overlap-conc}.
\end{proof}

\begin{proposition}[Base case at $t = 0$]\label{prop:sg-tal:base-case}
  \lean{ProbabilityTheory.cosh_add_mul_sinh_mul_exp_neg_le,
        ProbabilityTheory.integral_exp_mul_sq_gaussianReal,
        ProbabilityTheory.integrable_exp_mul_sq_gaussianReal,
        SpinGlasses.integral_exp_sq_le_of_mgf_le,
        SpinGlasses.sum_sum_exp_overlap_site,
        SpinGlasses.sum_sum_decoupled_overlap_exp,
        SpinGlasses.integral_decoupled_pair_exp,
        SpinGlasses.decoupled_field_eq,
        SpinGlasses.integral_decoupled_pair_exp_le,
        SpinGlasses.integral_decoupled_exp_overlap_sq_le}
  \leanok
Let $N\ge1$, $\beta,h\in\mathbb R$, and let $q$ satisfy the fixed-point equation $q = \mathbb{E}\tanh^2(\beta\sqrt q\,z+h)$ (so $q\in[0,1]$).
For the decoupled model (non-interacting spins with fields $\beta\sqrt q\, z_i + h$, two independent replicas) and every $0 \le \lambda < 1/2$ (Proposition~1.4.8):
\[
\mathbb{E}\big\langle \exp\big(\lambda N (R_{1,2} - q)^2\big)\big\rangle \le \frac{1}{\sqrt{1 - 2\lambda}} .
\]
\end{proposition}
\begin{proof}
\leanok
\emph{Step 1 (one site).} For $\varepsilon=\pm1$ with $\mathbb{E}\varepsilon=q$, Hoeffding's lemma (Mathlib, applied to the explicit two-point law) gives $(\cosh u + q\sinh u)e^{-qu}=\mathbb{E}e^{u(\varepsilon-q)}\le e^{u^2/2}$.
\emph{Step 2 (replicas factorize).} For fixed fields $Y_i$, summing first over $\sigma^2$ and then over $\sigma^1$ with the hypercube factorization gives
$\langle\exp(u\sum_i(\sigma^1_i\sigma^2_i-q))\rangle=\prod_i e^{-uq}(\cosh u+\tanh^2(Y_i)\sinh u)$, because $\sum_{s,t=\pm1}e^{(s+t)y+u(st-q)}=4\cosh^2 y\;e^{-uq}(\cosh u+\tanh^2 y\sinh u)$.
\emph{Step 3 (disorder).} The coordinates $z_i$ are independent, so the expectation over $z$ is the $N$-th power of $e^{-uq}(\cosh u+\mathbb{E}\tanh^2(Y)\sinh u)=e^{-uq}(\cosh u+q\sinh u)$, which is at most $e^{Nu^2/2}$ by Step~1 (Eq.~(1.105)).
\emph{Step 4 (Gaussian linearization).} For $X=N(R_{1,2}-q)$ one has $\mathbb{E}\langle e^{uX}\rangle\le e^{Nu^2/2}$ for all $u$; writing $e^{\theta X^2/2}=\mathbb{E}_g e^{\sqrt\theta\, g X}$ with $g\sim\mathcal N(0,1)$, exchanging expectations (Fubini; $X$ is bounded) and using the bound with $u=\sqrt\theta g$ gives $\mathbb{E}\langle e^{\theta X^2/2}\rangle\le\mathbb{E}_g e^{N\theta g^2/2}=(1-N\theta)^{-1/2}$.
With $\theta=2\lambda/N$ this is the claim.
\emph{Formalization note.} The fixed-point equation is a hypothesis; only existence of a solution is formalized (Theorem~\ref{thm:sg-tal:fixed-point}).
Formalized in \texttt{SpinGlasses/SKLatalaBase.lean}.
\end{proof}

\begin{theorem}[Exponential overlap concentration]\label{thm:sg-tal:overlap-conc}
  \lean{SpinGlasses.sk_overlap_concentration_exp,
        SpinGlasses.sk_path_expected_sq_overlap_le}
  \leanok
Let $N\ge1$, $\beta,h\in\mathbb R$ and let $q=\mathbb{E}\tanh^2(\beta\sqrt q\,z+h)$.
For $s\ge0$ with $2s+4\beta^2<1$ (Theorem~1.4.1):
\[
\nu\left(\exp s N (R_{1,2} - q)^2\right) \le \frac{1}{\sqrt{1 - 2s - 4\beta^2}} .
\]
Moreover, if $\beta^2<1/4$ then for every $t\in[0,\pi/2]$ the expected squared overlap under the smart-path energy at time $t$ satisfies (Eqs.~(1.89) and (1.103))
\[
\nu_t\big((R_{1,2}-q)^2\big)\le \frac{K}{N},\qquad K=\frac{(1/2-2\beta^2)^{-1/2}-1}{1/4-\beta^2}.
\]
\end{theorem}
\begin{proof}
\leanok
Apply Corollary~\ref{cor:sg-tal:compensated-mono} with $\lambda=s+2\beta^2<1/2$ and then Proposition~\ref{prop:sg-tal:base-case}:
$\nu(\exp sN(R_{1,2}-q)^2)\le\nu_{\mathrm{dec}}(\exp\lambda N(R_{1,2}-q)^2)\le(1-2\lambda)^{-1/2}=(1-2s-4\beta^2)^{-1/2}$.
For the second claim take $\lambda=1/4+\beta^2$ and $t\in[0,\pi/2]$.
The compensated inequality at time $t$ gives $\nu_t(\exp(aN(R_{1,2}-q)^2))\le(1-2\lambda)^{-1/2}$ with $a=\lambda-2\beta^2\cos^2t\ge 1/4-\beta^2$.
From $e^x\ge1+x$ we get $1+aN\,\nu_t((R_{1,2}-q)^2)\le(1-2\lambda)^{-1/2}$, hence $\nu_t((R_{1,2}-q)^2)\le K/N$.
\emph{Formalization notes.} (i) The fixed-point equation $q=\mathbb{E}\tanh^2(\beta\sqrt q\,z+h)$ is a hypothesis of both statements; only the existence of a solution is proved (Theorem~\ref{thm:sg-tal:fixed-point}).
(ii) The book's assumption $\beta<1/2$ is stated as $\beta^2<1/4$ (in the first claim it is implied by $2s+4\beta^2<1$ with $s\ge0$).
(iii) Eq.~(1.88), the bound $(Kk/N)^k$ on all even moments, is not formalized; only the exponential bound and the second moment are.
Formalized in \texttt{SpinGlasses/SKOverlapConcentration.lean}.
\end{proof}

\begin{theorem}[The Aizenman--Lebowitz--Ruelle high-temperature free energy formula]\label{thm:sg-tal:alr}
  \lean{SpinGlasses.ge_endpoint_of_trig_square_deriv,
        SpinGlasses.exists_sk_smart_path_profile,
        SpinGlasses.sk_free_energy_sub_rs_le,
        SpinGlasses.tendsto_sk_free_energy_of_lt}
  \leanok
For $\beta^2<1/4$ and $q=\mathbb{E}\tanh^2(\beta\sqrt q\,z+h)$ (Theorem~1.4.10):
\[
\Big|p_N(\beta, h) - \Big(\log2+\mathbb{E}\log\cosh(\beta z\sqrt q+h)+\tfrac{\beta^2}{4}(1-q)^2\Big)\Big| \le \frac{\beta^2}{4}\cdot\frac{(1/2-2\beta^2)^{-1/2}-1}{(1/4-\beta^2)\,N}.
\]
Consequently $p_N(\beta,h)\to\operatorname{SK}(\beta,h)=\log2+\mathbb{E}\log\cosh(\beta z\sqrt q+h)+\frac{\beta^2}{4}(1-q)^2$.
\end{theorem}
\begin{proof}
\leanok
Let $\varphi(s)=N^{-1}\mathbb E\log Z_s$ along the trigonometric smart path.
By Theorem~\ref{thm:sg-tal:guerra-rs-bound} (profile lemma) $\varphi'(s)=\frac{\beta^2}{2}\sin s\cos s\,(A(s)-(1-q)^2)$ with $A(s)=\nu_s((R_{1,2}-q)^2)\ge0$, $\varphi(0)=p_N$ and $\varphi(\pi/2)=\log 2+\mathbb E\log\cosh(\beta z\sqrt q+h)$.
Since $0\le A(s)\le K'/N$ on $[0,\pi/2]$ by Theorem~\ref{thm:sg-tal:overlap-conc}, the functions $\varphi(s)+\frac{\beta^2}{4}(1-q)^2\sin^2 s$ (nondecreasing) and $\varphi(s)+\frac{\beta^2}{4}((1-q)^2-K'/N)\sin^2 s$ (nonincreasing) give $\varphi(0)\le\varphi(\pi/2)+\frac{\beta^2}{4}(1-q)^2$ and $\varphi(0)\ge\varphi(\pi/2)+\frac{\beta^2}{4}((1-q)^2-K'/N)$, i.e.\ the claim.
\emph{Formalization note.} This is done by monotonicity instead of integrating the derivative (same bound, no integrability to check).
As in Theorem~\ref{thm:sg-tal:overlap-conc}, $q$ is a hypothesis (fixed-point equation) and $\beta<1/2$ is stated as $\beta^2<1/4$.
\emph{Not formalized:} the identification of the solution $q$ with the minimizer of Guerra's bound and the uniqueness of the fixed point for $\beta<1$; the thermodynamic limit of the Guerra--Toninelli theorem is stated there for the full-matrix disorder without field, while the limit here is derived directly from the $O(1/N)$ bound for the pair disorder with field.
Formalized in \texttt{SpinGlasses/SKAizenmanLebowitzRuelle.lean}.
\end{proof}

\section{Spin Decorrelation, Asymptotic Product Measures & Projections (Talagrand Sections~1.4--1.5)}\label{sec:sg-tal:decorrelation}

The final component of the high-temperature theory connects overlap concentration to local spin independence and geometric projections.

\begin{proposition}[Spin decorrelation from overlap concentration]\label{prop:sg-tal:spin-decorr}
For any $p \ge 1$, site symmetry and overlap concentration imply truncated spin correlations vanish (Proposition~1.4.14, Eq.~(1.132)):
\[
\mathbb{E} \left(\langle \sigma_1 \cdots \sigma_p \rangle - \langle \sigma_1 \rangle \cdots \langle \sigma_p \rangle\right)^2 \le K(p) \nu\left((R_{1,2} - q)^2\right) \le \frac{K(p)}{N}.
\]
\end{proposition}
\begin{proof}
Center the spins $\dot{\sigma}_i = \sigma_i - \langle \sigma_i \rangle$.
By site permutation symmetry, the sum over all $N(N-1)\cdots(N-p+1)$ disjoint $p$-tuples of spins of $\langle \dot{\sigma}_{i_1} \dots \dot{\sigma}_{i_p} \rangle^2$ is bounded by $N^p \nu((\frac{\dot{\sigma}^1 \cdot \dot{\sigma}^2}{N})^p) \le K(p) N^p \nu((R_{1,2}-q)^2)$.
Dividing by $N^p$ establishes the result.
\end{proof}

\begin{theorem}[Gibbs measure is asymptotically product on finite marginals]\label{thm:sg-tal:asymp-product}
For any fixed $p \in \mathbb{N}$, as $N \to \infty$ the marginal distribution of the first $p$ spins $(\sigma_1, \dots, \sigma_p)$ under the Gibbs measure converges in total variation distance to the product measure $\bigotimes_{i=1}^p \mu_i$, where $\mu_i(\sigma_i) = \frac{1 + \sigma_i \langle \sigma_i \rangle}{2}$ (Theorem~1.4.15).
\end{theorem}
\begin{proof}
Any function on $\{-1, 1\}^p$ expands as a linear combination of spin monomials $\sigma_S = \prod_{i \in S} \sigma_i$ for $S \subseteq \{1, \dots, p\}$.
By Proposition~\ref{prop:sg-tal:spin-decorr}, $\langle \sigma_S \rangle - \prod_{i \in S} \langle \sigma_i \rangle \to 0$ in $L^2$, ensuring total variation convergence.
\end{proof}

\begin{theorem}[High-dimensional random projections under Gibbs measure]\label{thm:sg-tal:random-projection}
Let $\mathbf{g} = (g_1, \dots, g_N)$ be standard Gaussian random variables independent of the disorder.
Under the high-temperature Gibbs measure, the random projection $\frac{\mathbf{g} \cdot \sigma}{\sqrt{N}}$ converges to a Gaussian mixture with variance $1 - q$ centered at the random barycenter $\frac{\mathbf{g} \cdot \langle \sigma \rangle}{\sqrt{N}}$:
\[
\left\langle \exp\left(t \frac{\mathbf{g} \cdot (\sigma - \langle \sigma \rangle)}{\sqrt{N}}\right) \right\rangle \xrightarrow[N \to \infty]{\mathbb{P}} \exp\left(\frac{t^2}{2}(1 - q)\right).
\]
\end{theorem}
\begin{proof}
Compute the conditional variance $\operatorname{Var}_\sigma(\frac{\mathbf{g} \cdot \sigma}{\sqrt{N}}) = \frac{1}{N}\sum_{i,j} g_i g_j (\langle \sigma_i \sigma_j \rangle - \langle \sigma_i \rangle \langle \sigma_j \rangle)$.
By spin decorrelation, off-diagonal terms vanish, and the diagonal contributes $\frac{1}{N}\sum_i g_i^2 (1 - \langle \sigma_i \rangle^2) \to 1 - q$ by the law of large numbers.
\end{proof}
